\documentclass[11pt]{article}
\usepackage[margin=1in]{geometry}
\usepackage[T1]{fontenc}
\usepackage{lmodern,amsmath,amssymb,amsthm,mathtools,hyperref}
\setlength{\emergencystretch}{3em}
\newtheorem{theorem}{Theorem}[section]
\newtheorem{lemma}[theorem]{Lemma}
\newtheorem{corollary}[theorem]{Corollary}
\newtheorem{proposition}[theorem]{Proposition}
\newcommand{\E}{\mathbb E}
\newcommand{\tr}{\operatorname{tr}}
\title{Four-phase Hadamard constructions\\and a quantitative Fourier obstruction}
\author{Joshua Mark\thanks{ORCID: \href{https://orcid.org/0009-0003-5905-2382}{0009-0003-5905-2382}. AI tools assisted with derivation, independent review, computation and exposition.}}
\date{20 November 2026 --- research draft; priority not established}
\begin{document}\maketitle
\section{A four-phase route and a quantitative Fourier obstruction}
\label{sec:four-phase-route}

A quaternary Hadamard matrix $B$ of order $m$, also denoted
$\mathrm{BH}(m,4)$, has entries in $\mu_4=\{1,-1,i,-i\}$ and
$BB^*=mI_m$. Such matrices give real Hadamard matrices of order
$2m$. This is a classical construction, not a new existence theorem;
see Butson, \emph{Generalized Hadamard matrices} (1962),
\url{https://doi.org/10.1090/S0002-9939-1962-0142557-0},
and Egan et al., the morphism discussion in
\url{https://ajc.maths.uq.edu.au/pdf/73/ajc_v73_p501.pdf}.

This route permits arbitrary complex phases as an initial search
space, which is nonempty in every dimension, but the final four-phase
constraint must actually be met. The results below give an exact
phase criterion and a quantitative obstruction to some obvious
Fourier deformations. Priority for the combined observations remains
unestablished; neither gives an all-order existence theorem.

\subsection{Exact realification and the missing monomial structure}
For a complex matrix $Z$ write $\mathcal R(Z)$ for its realification,
replacing $x+iy$ by $\left(\begin{smallmatrix}x&-y\\y&x\end{smallmatrix}\right)$.
It respects products and sends conjugate transpose to real transpose.
For $B\in\mathrm{BH}(m,4)$, put
\[
                   H=\mathcal R((1+i)B).
\]
Its entries are signs and $HH^T=2mI_{2m}$. Thus proving
$\mathrm{BH}(2t,4)$ existence for every positive $t$ would suffice
for the full real Hadamard conjecture. No converse or proof of this
stronger target is being assumed.

\begin{proposition}[A signed complex-structure criterion]
\label{quat:structure}
A real Hadamard matrix $H$ of order $2m$ is equivalent, by row/column
permutations and sign changes, to such a realification if and only if there are signed
permutation matrices $J_{\rm row},J_{\rm col}$ with
\[
 J_{\rm row}^2=J_{\rm col}^2=-I,\qquad
                         HJ_{\rm col}=J_{\rm row}H.
\]
\end{proposition}
\begin{proof}
Each signed permutation squaring to $-I$ pairs coordinates, with
opposite signs on each two-cycle. Signed permutations of coordinates
put each into the canonical block form
$J=\operatorname{diag}(\left(\begin{smallmatrix}0&-1\\1&0\end{smallmatrix}\right))$.
Transform rows and columns independently. The intertwining equation
becomes $H'J=JH'$, so every two-square block is
$\left(\begin{smallmatrix}a&-b\\b&a\end{smallmatrix}\right)$ with
$a,b\in\{\pm1\}$. Hence $H'=\mathcal R(Z)$ and
$ZZ^*=2mI_m$. Dividing $Z$ by $1+i$ gives a quaternary matrix
of Gram $mI_m$. The converse follows by taking the canonical $J$
on both sides of the displayed realification.
\end{proof}

The signed-permutation condition is essential. For any even-order
scaled orthogonal $H$ and any orthogonal complex structure $J_{\rm col}$,
one can set $J_{\rm row}=HJ_{\rm col}H^T/(2m)$. This always gives
an intertwiner and a real orthogonal complex structure, but usually
not a signed permutation. It need not pair individual coordinates,
so it does not give a four-phase matrix. The checker verifies this
failure for the default adjacent pairing on a normalized Paley $H_{12}$;
it does not claim that no other pairing works.

\subsection{An invariant that detects four-phase accessibility}
Let $F$ be any complex Hadamard matrix of order $n$ with unit-modulus
entries. Write $F^{\circ4}$ for its entrywise fourth power and define
\[
                 \rho_4(F)=\frac{\|F^{\circ4}\|_{\rm op}}n.
\]
Multiplying rows or columns by unit phases, or permuting them, does
not change $\rho_4$, since these operations induce unitary monomial
operations on $F^{\circ4}$.

\begin{theorem}[Exact criterion and quantitative separation]
\label{quat:phase-defect}
One has $\rho_4(F)\le1$, with equality if and only if $F$ can be
converted into a quaternary Hadamard matrix by row and column phases.
For every matrix $F'$ equivalent to $F$ under row/column phases and
permutations, and every matrix $Q$ with entries in $\mu_4$,
\[
              \frac{\|F'-Q\|_F^2}{n^2}
                    \ge\frac{1-\rho_4(F)}8.
\]
The comparison matrix $Q$ need not be Hadamard.
\end{theorem}
\begin{proof}
The Frobenius norm of $M=F^{\circ4}$ is $n$, so its operator norm
is at most $n$. Equality is equivalent to $M$ having rank one.
A rank-one matrix with unit-modulus entries factors as $M_{ij}=u_i v_j$
with unit $u_i,v_j$; choose fourth-root row and column phases to
remove them. The transformed entries then have fourth power one,
and the Hadamard Gram is preserved. The reverse implication is immediate.

For the quantitative bound, set $M'=(F')^{\circ4}$ and let $J_n$
be the all-one matrix. Since $|z^4-w^4|\le4|z-w|$ on the unit circle,
\[
 16\|F'-Q\|_F^2\ge\|M'-J_n\|_F^2
 =2n^2-2\operatorname{Re}(\mathbf1^TM'\mathbf1)
 \ge2n^2-2n\|M'\|_{\rm op}.
\]
Use invariance of the operator norm under the phase and permutation
changes to obtain the result.
\end{proof}

This criterion is a spectral version of the classical four-cycle
phase invariants of complex Hadamard equivalence. The relevant
background and the block deformation construction below are discussed
in Tadej--\.Zyczkowski, \emph{A concise guide to complex Hadamard matrices},
\url{https://arxiv.org/abs/quant-ph/0512154}.
It does not show that an optimizer of $\rho_4$ reaches one.

\subsection{Fourier blocks retain a definite obstruction}
Let $\mathcal F_m=(\zeta_m^{ij})_{0\le i,j<m}$ and $d=\gcd(m,4)$.
Direct character orthogonality gives
\[
 (\mathcal F_m^{\circ4}(\mathcal F_m^{\circ4})^*)_{ij}
 =\begin{cases}m&m\mid4(i-j),\\0&\text{otherwise}.\end{cases}
\]
There are $m/d$ all-one blocks of size $d$ after row reordering, so
\[
 \operatorname{rank}(\mathcal F_m^{\circ4})=m/d,\qquad
 \rho_4(\mathcal F_m)=\sqrt{d/m}.
\]
In particular Fourier row/column phases can produce a quaternary
matrix exactly when $m$ divides four. Arbitrary phases do not remove
this obstruction.

More generally take an order-$r$ complex Hadamard matrix $A$,
unit phases $u_{bi}$, and the classical block family
\[
 F_{(a,i),(b,j)}=A_{ab}u_{bi}\zeta_m^{ij}.
\]
It is Hadamard: inner products vanish in the Fourier factor for
distinct $i$, and otherwise reduce to those of $A$.
Every fourth-power block has operator norm $\sqrt{md}$.
For a vector partitioned into $r$ blocks, the triangle and
Cauchy--Schwarz inequalities therefore bound the full operator norm
by $r\sqrt{md}$. Consequently
\[
 \rho_4(F)\le\sqrt{d/m},\qquad
 \inf_{F'\sim F,\ Q\in\mu_4^{rm\times rm}}
 \frac{\|F'-Q\|_F^2}{(rm)^2}
 \ge\frac{1-\sqrt{d/m}}8.
\]
Thus these Fourier-block deformations cannot reach the four-phase
target when $m\nmid4$, even with arbitrarily chosen deformation phases.
This excludes one search family, not every continuous path through
complex Hadamard matrices.

At order six the distinction is concrete: Fourier $\mathcal F_6$
has $\rho_4=1/\sqrt3<1$, whereas a symmetric conference matrix $C$
of order six gives $B=I+iC\in\mathrm{BH}(6,4)$, with $\rho_4(B)=1$.
Its realification gives the classical $H_{12}$. The checker verifies
this and the corresponding orders fourteen and eighteen in exact
Gaussian-integer arithmetic, producing real orders twenty-eight
and thirty-six. The conference construction comes from Paley's
known prime-field matrices, not an unknown-order existence argument.

The sufficient research target is now precise: find a four-phase
complex Hadamard matrix at every even order, or force $\rho_4=1$
somewhere in each corresponding complex-Hadamard space.
Those spaces are nonempty, but their Fourier phase or block
families alone cannot supply the general conclusion.

\end{document}
